%%%%%%%%%%%%%%%%%%%%%%%%
% $Autor: MansourAhmadyPhoulady, RanjeetsingTawar, HemanthPrabhakaran$
% $Pfad: Car_Licence_Plate_Identification_with_a_JetsonNano$
% $Version: 4.0.0 $
%
% !TeX encoding = utf8
%
%%%%%%%%%%%%%%%%%%%%%%%%

\chapter{Program Flowchart}

\section{Flowchart}

\begin{itemize}
	\item Start - Initiation of the process.
	\item License Plate Detection - The webcam is active and searches for a license plate.
	\item License Plate Verification - Determines if the license plate is recognized as valid.
	\item Retry Detection - If the license plate is not valid, the system retries the detection.
	\item Highlight Plate - Adds a rectangular highlight around the identified plate.
	\item Extract Plate Text - Captures and processes the text from the license plate.
	\item Verify Accuracy - Assesses if the extracted text meets quality and accuracy standards.
	\item Accuracy Check - Determines if the accuracy is satisfactory.
	\item Data Omission - If the data is not accurate enough, it is not stored.
	\item Data Storage - Accurate data is saved in a CSV file, each entry tagged with a unique identifier (UUID).
	\item End - The process concludes.
\end{itemize}

\begin{tikzpicture}[node distance=0.5cm and 1cm, auto,
	block/.style={
		rectangle, 
		draw, 
		fill=blue!20, 
		text width=8em, 
		align=center, 
		rounded corners, 
		minimum height=4em
	},
	cloud/.style={
		draw, 
		ellipse,
		fill=green!20, 
		minimum height=1em
	},
	decision/.style={
		diamond, 
		aspect=2, 
		draw, 
		fill=yellow!20, 
		text width=7em, 
		align=center, 
		inner sep=-2pt
	},
	line/.style={
		draw, 
		-latex
	},
	every node/.style={font=\sffamily\small}
	]
	
	% Nodes
	\node [cloud] (init) {Begin};
	\node [block, below=of init] (detect) {License Plate Detection};
	\node [decision, below=of detect] (verify) {License Plate Verification};
	\node [block, left=of verify, xshift=-2cm] (retry) {Retry Detection};
	\node [block, below=of verify] (highlight) {Highlight Plate};
	\node [block, below=of highlight] (extract) {Extract Plate Text};
	\node [decision, below=of extract] (accuracyCheck) {Verify Accuracy};
	\node [cloud, left=of accuracyCheck, xshift=-2cm] (omit) {Data Omission};
	\node [block, below=of accuracyCheck, node distance=2cm] (store) {Data Storage};
	\node [cloud, below=of store] (finish) {End};
	
	% Paths
	\path [line] (init) -- (detect);
	\path [line] (detect) -- (verify);
	\path [line] (verify) -- node [near start] {yes} (highlight);
	\path [line] (verify) -- node [near start] {no} (retry);
	\path [line] (retry) |- (detect);
	\path [line] (highlight) -- (extract);
	\path [line] (extract) -- (accuracyCheck);
	\path [line] (accuracyCheck) -- node [near start] {yes} (store);
	\path [line] (accuracyCheck) -- node [near start] {no} (omit);
	\path [line] (store) -- (finish);
\end{tikzpicture}

\section{Description}

\begin{itemize}
	\item \textbf{Start - Initiation of the process.} This marks the beginning of the automated system operation. It initializes the necessary software components and hardware interfaces for the license plate recognition process.
	
	\item \textbf{License Plate Detection -} The system activates the connected webcam and begins a real-time search for license plates within the camera's field of view. This step involves image capture and the application of machine learning or computer vision algorithms to identify potential license plate regions.
	
	\item \textbf{License Plate Verification -} Once a license plate is detected, the system checks its format and characteristics against a predefined set of criteria to determine its validity. This step ensures that the detected object is indeed a license plate and not a false positive.
	
	\item \textbf{Retry Detection -} If the license plate is deemed invalid or not recognized, the system automatically retries the detection process. This loop continues until a valid license plate is detected, ensuring the system's robustness against detection errors.
	
	\item \textbf{Highlight Plate -} After successful verification, the system visually highlights the identified license plate within the image. This usually involves drawing a rectangular border around the license plate, making it easier for operators to verify the detection visually.
	
	\item \textbf{Extract Plate Text -} With the license plate highlighted, the system proceeds to capture and process the text displayed on the license plate. Advanced optical character recognition (OCR) algorithms are employed to read and interpret the alphanumeric characters accurately.
	
	\item \textbf{Verify Accuracy -} Following text extraction, the system assesses the quality and accuracy of the captured text. This involves comparing the extracted text against expected formats and performing error checking to minimize the likelihood of inaccuracies.
	
	\item \textbf{Accuracy Check -} An additional step to ensure the extracted text meets the system's quality standards. If the accuracy is deemed satisfactory, the process moves forward; otherwise, corrective measures or manual checks might be initiated.
	
	\item \textbf{Data Omission -} In cases where the data does not meet the accuracy or quality standards, it is omitted from further processing. This decision helps maintain the integrity of stored data, ensuring only reliable information is saved.
	
	\item \textbf{Data Storage -} Accurate and verified license plate data is saved in a structured CSV file. Each entry is tagged with a unique identifier (UUID) to facilitate easy retrieval and reference. This step is crucial for maintaining a record of detected and processed license plates.
	
	\item \textbf{End -} This final step marks the conclusion of the process. The system may either shut down or become ready to initiate a new cycle, depending on the operational settings and requirements.
\end{itemize}


\section{Actual state of the working project}
The Automated License Plate Recognition (ALPR) system is a sophisticated integration of computer vision, machine learning, and data management technologies, designed to detect, recognize, and process vehicle license plates in real-time. This system finds its applications in traffic monitoring, parking management, access control, and law enforcement, aiming to automate the identification and logging of vehicle license plates efficiently and accurately.

\subsection{Initialization}
The ALPR system initiates operations by preparing its hardware components (cameras, webcams) and software modules. This stage involves loading machine learning models for image processing and configuring the database for subsequent data storage.

\subsection{License Plate Detection}
Utilizing live camera feeds, the system applies computer vision algorithms to identify objects resembling license plates. It leverages advanced image analysis to distinguish license plates from other objects under various environmental conditions.

\subsection{Verification and Highlighting}
Detected objects are verified against predefined criteria to ensure authenticity as license plates. Verified plates are then highlighted within the image for visual confirmation and further processing.

\subsection{Text Extraction and Accuracy Verification}
Through Optical Character Recognition (OCR) technology, the system extracts text from the highlighted license plates. Extracted text undergoes a rigorous accuracy verification process to ensure it meets the set quality standards, critical for maintaining data reliability.

\subsection{Data Management}
Accurate data is processed for storage in a structured CSV format, with each entry uniquely identified for easy access and reference. Inaccurate data is omitted to maintain the integrity of stored information.

\subsection{Continuous Operation and Error Handling}
Designed for continuous operation, the system includes error-handling mechanisms to manage detection retries or manual intervention prompts, ensuring uninterrupted functionality.

\subsection{Conclusion}
The process concludes with the successful storage of data or manual termination. The system's design allows for automatic cycling for continuous operation, showcasing its adaptability and efficiency.
